\chapter{Vector Spaces, Superposition, and the First Quantum Postulates}

We ended the previous lecture with a question about the uncertainty
principle. Before building any new formalism, let us settle it. The apparent
loophole is instructive because it distinguishes the momentum a particle had
before a measurement from the momentum it has after the measurement.

\section{Which Velocity Have We Measured?}

Suppose a first position measurement locates a particle only within a region
of width \(\delta\). Wait a time \(T\), make a second measurement with the
same resolution, and let the centers of the two regions be separated by
\(L\). The inferred average velocity is
\begin{equation}
 v_{\mathrm{avg}}\simeq \frac{L}{T},
 \qquad
 \Delta v_{\mathrm{avg}}\sim \frac{\delta}{T}.
 \label{eq:l02-coarse-velocity}
\end{equation}
Numerical factors in the second estimate are unimportant. By making \(T\)
large, we can make \(\delta/T\) as small as we please. Have we measured
position to accuracy \(\delta\) and velocity almost exactly?

\begin{figure}[tbp]
\centering
\includegraphics[width=0.68\textwidth]{lecture_02_figure_01.png}
\caption{Two coarse position readings separated by a long time. The board
displays the displacement \(L\), the resolution scale \(\delta\), and the
velocity uncertainty of order \(\delta/T\).}
\lectureframe{00:03:30}
\end{figure}

\begin{classroomqa}
\audiencequestion{Can two imprecise position measurements, separated by a
very long time, evade the uncertainty principle?
\lecturetimestamp{00:00:17--00:04:17}}
\lecturerresponse{No. Equation~\eqref{eq:l02-coarse-velocity} gives the
velocity after the first position measurement. To localize the particle on
the scale \(\delta\), the probe must contain wavelengths no larger than that
scale. A photon capable of doing so carries momentum of order
\(\hbar/\delta\), and its scattering gives the particle a kick of the same
order. The later trajectory therefore reveals the momentum of the disturbed
state, not the momentum immediately before localization.}
\end{classroomqa}

Trying to verify the result does not restore the missing information.
Repeating the procedure changes the momentum again by an amount of order
\(\hbar/\delta\). Classical physics permits an ideal limit in which position
is observed ever more gently while the velocity is left alone. Quantum
mechanics does not. This is the point of the uncertainty principle, and with
it settled we can begin the mathematics needed for states.

\section{The One-Dimensional Model: Complex Numbers}

The next subject is dry, but it is not difficult. We need a quick review of
complex numbers only to establish notation. Write
\begin{equation}
 z=x+iy,
 \qquad
 z^*=x-iy,
 \qquad
 i^2=-1.
 \label{eq:l02-complex-number}
\end{equation}
In the complex plane, conjugation reflects a point across the real axis.
Complex numbers can be added and multiplied by arbitrary complex scalars, so
\(\mathbb C\) is itself a one-dimensional complex vector space. It is the
simplest model of the structures we shall use.

In this one-dimensional model, regard \(z\) as a ket and \(z^*\) as its
corresponding bra. Complex conjugation exchanges the two. If
\[
 z_1=x_1+iy_1,
 \qquad
 z_2=x_2+iy_2,
\]
then
\begin{align}
 z_1z_2
 &= (x_1x_2-y_1y_2)+i(x_1y_2+y_1x_2),\\
 (z_1z_2)^*
 &=z_1^*z_2^*.
 \label{eq:l02-conjugate-product}
\end{align}
For scalar numbers the factors commute. It is nevertheless useful to write
the related identity
\begin{equation}
 (z_1^*z_2)^*=z_2^*z_1,
 \label{eq:l02-reversed-product}
\end{equation}
because the reversal of bra and ket will survive when the intervening
objects no longer commute.

\begin{figure}[tbp]
\centering
\includegraphics[width=0.66\textwidth]{lecture_02_figure_02.png}
\caption{Complex conjugation of a product on the board. The corrected
relation is given cleanly in Eqs.~\eqref{eq:l02-conjugate-product} and
\eqref{eq:l02-reversed-product}.}
\lectureframe{00:14:30}
\end{figure}

\begin{classroomqa}
\audiencequestion{Can the complex number be written in exponential form?
\lecturetimestamp{00:15:12--00:15:18}}
\lecturerresponse{Certainly. One may write \(z=re^{i\theta}\), but we do not
need that form yet. The Cartesian form makes conjugation and the coming
vector-space rules visible with less machinery.}
\end{classroomqa}

\section{What Makes a Complex Vector Space?}

A complex vector space is a collection of objects that may be added and
multiplied by complex numbers without leaving the collection. If
\(|A\rangle\) and \(|B\rangle\) are vectors and
\(\alpha,\beta\in\mathbb C\), then
\begin{equation}
 \alpha|A\rangle+\beta|B\rangle
 \label{eq:l02-linear-combination}
\end{equation}
must be another vector in the same space. The vectors need not be arrows.
What matters is the algebra of addition and scalar multiplication.

Several near-examples make the definition precise. Ordinary arrows form a
real vector space, but not a complex vector space unless an additional rule
for multiplication by \(i\) is supplied. The real numbers form a real vector
space, not a complex one. The positive real numbers do not even form a real
vector space: multiplying a positive number by \(-1\) leaves the set.

\subsection{Functions and Columns}

Complex-valued functions of a real variable form a complex vector space. If
\(\psi(x)\) and \(\phi(x)\) are such functions, then
\begin{equation}
 (\alpha\psi+\beta\phi)(x)
 =\alpha\psi(x)+\beta\phi(x)
 \label{eq:l02-function-combination}
\end{equation}
is another complex-valued function.

A column vector is the same idea with a discrete argument. For four
components,
\begin{equation}
 |a\rangle=
 \begin{pmatrix}a_1\\a_2\\a_3\\a_4\end{pmatrix},
 \qquad
 |b\rangle=
 \begin{pmatrix}b_1\\b_2\\b_3\\b_4\end{pmatrix},
\end{equation}
and
\begin{equation}
 \alpha|a\rangle+\beta|b\rangle
 =
 \begin{pmatrix}
  \alpha a_1+\beta b_1\\
  \alpha a_2+\beta b_2\\
  \alpha a_3+\beta b_3\\
  \alpha a_4+\beta b_4
 \end{pmatrix}.
 \label{eq:l02-column-combination}
\end{equation}
The index \(1,2,3,4\) plays the role of a discrete coordinate.

\begin{figure}[tbp]
\centering
\includegraphics[width=0.70\textwidth]{lecture_02_figure_03.png}
\caption{Componentwise addition of two column vectors. A column may be read
as a function of a discrete index.}
\lectureframe{00:24:30}
\end{figure}

\subsection{Dimension and Spanning}

The dimension of a vector space is the minimum number of independent vectors
needed to span it. One nonzero vector spans a line. Two noncollinear vectors
span a plane: every vector \(|v\rangle\) in that plane can be written
\begin{equation}
 |v\rangle=r_1|v_1\rangle+r_2|v_2\rangle.
 \label{eq:l02-plane-span}
\end{equation}
A third vector in the same plane adds no new direction, because it is already
a linear combination of the first two. In three dimensions we need three
noncoplanar vectors.

\begin{figure}[tbp]
\centering
\includegraphics[width=0.67\textwidth]{lecture_02_figure_04.png}
\caption{Two noncollinear vectors spanning a plane. The expression at left is
partly obscured, while the geometric pair is clear.}
\lectureframe{00:28:59}
\end{figure}

\begin{figure}[tbp]
\centering
\begin{tikzpicture}[>=Latex,scale=1.0]
  \coordinate (O) at (0,0);
  \draw[->,thick] (O) -- (4.2,0.5) node[right] {$|v_2\rangle$};
  \draw[->,thick] (O) -- (2.0,3.0) node[above] {$|v_1\rangle$};
  \draw[->,very thick,blue!65!black] (O) -- (3.0,1.9)
    node[above right] {$|v\rangle$};
  \draw[dashed,gray] (2.0,3.0) -- (5.0,3.5);
  \draw[dashed,gray] (4.2,0.5) -- (5.0,3.5);
\end{tikzpicture}
\caption{A clean reconstruction of the spanning argument:
\(|v\rangle=r_1|v_1\rangle+r_2|v_2\rangle\).}
\end{figure}

Columns with four components form a vector space of dimension four. One
convenient basis is
\begin{equation}
 \begin{pmatrix}1\\0\\0\\0\end{pmatrix},
 \quad
 \begin{pmatrix}0\\1\\0\\0\end{pmatrix},
 \quad
 \begin{pmatrix}0\\0\\1\\0\end{pmatrix},
 \quad
 \begin{pmatrix}0\\0\\0\\1\end{pmatrix}.
 \label{eq:l02-standard-columns}
\end{equation}
Remove the last column and no linear combination can produce a nonzero
fourth component. By contrast, one-component columns are just complex
numbers and form a one-dimensional complex space.

The space of smooth complex functions has no finite spanning set; it is
infinite-dimensional. The details of bases in infinite-dimensional spaces
will not concern us yet. For now, the essential point is that functions and
columns are both concrete representations of abstract vectors:
\begin{equation}
 \psi(x)\longmapsto|\psi\rangle,
 \qquad
 \begin{pmatrix}\psi_1\\ \psi_2\\ \vdots\end{pmatrix}
 \longmapsto|\psi\rangle.
 \label{eq:l02-function-ket}
\end{equation}

\begin{figure}[tbp]
\centering
\includegraphics[width=0.69\textwidth]{lecture_02_figure_05.png}
\caption{The board changes viewpoint from the function \(\psi(x)\) to the
abstract ket \(|\psi\rangle\). The second line is still being written.}
\lectureframe{00:34:50}
\end{figure}

\section{Bras, Kets, and the Dual Space}

To every ket \(|A\rangle\) we associate a bra \(\langle A|\). A bra is a
linear functional: when it acts on a ket, it produces a complex number. Once
an inner product has been chosen, the Riesz correspondence identifies each
ket with a bra antilinearly. In finite-dimensional coordinates this is the
familiar conjugate transpose:
\begin{equation}
 |A\rangle=
 \begin{pmatrix}a_1\\a_2\\ \vdots\\a_d\end{pmatrix}
 \quad\longleftrightarrow\quad
 \langle A|=
 \begin{pmatrix}a_1^*&a_2^*&\cdots&a_d^*\end{pmatrix}.
 \label{eq:l02-ket-bra}
\end{equation}
For a function, the corresponding bra is represented by the conjugate
function \(\psi^*(x)\).

The correspondence is antilinear:
\begin{align}
 \alpha|A\rangle
 &\longleftrightarrow
 \langle A|\alpha^*,\\
 |A\rangle+|B\rangle
 &\longleftrightarrow
 \langle A|+\langle B|.
 \label{eq:l02-dual-rules}
\end{align}

\begin{figure}[tbp]
\centering
\includegraphics[width=0.66\textwidth]{lecture_02_figure_06.png}
\caption{Ket-to-bra correspondence and the conjugation of a scalar
coefficient.}
\lectureframe{00:41:00}
\end{figure}

\begin{classroomqa}
\audiencequestion{Is there any significance to writing \(\alpha^*\) on the
right of the bra?\lecturetimestamp{00:43:24--00:43:57}}
\lecturerresponse{Not while \(\alpha\) is an ordinary complex number,
because scalar multiplication commutes. It is partly a convention for
keeping formulas orderly. Position will matter when noncommuting operators
enter the calculation.}
\end{classroomqa}

\section{Inner Products}

For ordinary real vectors, the dot product is
\begin{equation}
 \mathbf a\mathbin{\boldsymbol\cdot}\mathbf b
 =|\mathbf a|\,|\mathbf b|\cos\theta
 =\sum_i a_i b_i.
 \label{eq:l02-real-dot}
\end{equation}
The inner product \(\langle A|B\rangle\) is its complex counterpart. We use
the convention standard in physics: it is linear in the ket and
conjugate-linear in the bra. Thus
\begin{align}
 \langle A|(\beta B+\gamma C)\rangle
 &=\beta\langle A|B\rangle+\gamma\langle A|C\rangle,\\
 \langle \alpha A|B\rangle
 &=\alpha^*\langle A|B\rangle.
 \label{eq:l02-sesquilinearity}
\end{align}
It also obeys conjugate symmetry,
\begin{equation}
 \langle A|B\rangle=\langle B|A\rangle^*.
 \label{eq:l02-conjugate-symmetry}
\end{equation}
Setting \(B=A\) shows that \(\langle A|A\rangle\) is real. Positivity is an
additional axiom:
\begin{equation}
 \langle A|A\rangle\geq 0,
 \qquad
 \langle A|A\rangle=0\ \Longleftrightarrow\ |A\rangle=0.
 \label{eq:l02-positive}
\end{equation}
It is not a consequence of conjugate symmetry alone. Relativity supplies a
useful warning: the Minkowski metric admits positive, negative, and null
norms, so it is not an inner product in this positive-definite sense.

For complex-valued functions on a real interval \(I\),
\begin{equation}
 \langle\phi|\psi\rangle
 =\int_I dx\,\phi^*(x)\psi(x).
 \label{eq:l02-function-inner-product}
\end{equation}

\begin{figure}[tbp]
\centering
\includegraphics[width=0.68\textwidth]{lecture_02_figure_07.png}
\caption{The function-space inner product being assembled on the board. The
complete formula is Eq.~\eqref{eq:l02-function-inner-product}.}
\lectureframe{00:56:30}
\end{figure}

For columns,
\begin{equation}
 |a\rangle=
 \begin{pmatrix}a_1\\a_2\\ \vdots\\a_d\end{pmatrix},
 \qquad
 \langle b|=
 \begin{pmatrix}b_1^*&b_2^*&\cdots&b_d^*\end{pmatrix},
\end{equation}
and therefore
\begin{equation}
 \langle b|a\rangle=\sum_{i=1}^{d}b_i^*a_i.
 \label{eq:l02-column-inner-product}
\end{equation}

\begin{classroomqa}
\audiencequestion{Are the coordinates of the vector \(b\) the starred
quantities?\lecturetimestamp{00:57:55--00:58:37}}
\lecturerresponse{The ket \(|b\rangle\) has coordinates \(b_i\). The
corresponding bra \(\langle b|\) has entries \(b_i^*\). The bra symbol already
indicates that conjugation has been made; one does not rename the ket
coordinates themselves.}
\end{classroomqa}

\begin{classroomqa}
\audiencequestion{Must the variable \(x\) in the function inner product be
real?\lecturetimestamp{00:58:59--00:59:42}}
\lecturerresponse{In the present construction, yes. We may use several real
variables, say \(x\) and \(y\), and integrate over each of them. Functions of
complex arguments can also be studied, but they are not needed for this
course at this point.}
\end{classroomqa}

\section{Orthonormal Bases and Components}

Choose \(d\) vectors \(\{|b_i\rangle\}_{i=1}^{d}\) in a
\(d\)-dimensional space. They are normalized when
\(\langle b_i|b_i\rangle=1\), with no sum on \(i\), and mutually orthogonal
when \(\langle b_i|b_j\rangle=0\) for \(i\neq j\). Both statements are
contained in
\begin{equation}
 \langle b_i|b_j\rangle=\delta_{ij},
 \qquad
 \delta_{ij}=
 \begin{cases}
  1,&i=j,\\
  0,&i\neq j.
 \end{cases}
 \label{eq:l02-orthonormal}
\end{equation}
An orthonormal set of \(d\) vectors in a \(d\)-dimensional space is a basis.
The columns in Eq.~\eqref{eq:l02-standard-columns} provide the standard
example. Their corresponding bras are rows. A blackboard slip between the
row and column forms was caught in the room; the corrected notation is the
one used here.

Any vector can be expanded in the basis:
\begin{equation}
 |v\rangle=\sum_{i=1}^{d}v_i|b_i\rangle.
 \label{eq:l02-basis-expansion}
\end{equation}
Multiply from the left by \(\langle b_j|\):
\begin{align}
 \langle b_j|v\rangle
 &=\sum_{i=1}^{d}v_i\langle b_j|b_i\rangle\\
 &=\sum_{i=1}^{d}v_i\delta_{ji}
 =v_j.
 \label{eq:l02-component-extraction}
\end{align}
The component \(v_j\) is the inner product with the \(j\)-th basis vector.
Substitution gives the form we shall use repeatedly,
\begin{equation}
 |v\rangle
 =\sum_{i=1}^{d}|b_i\rangle\langle b_i|v\rangle.
 \label{eq:l02-completeness-action}
\end{equation}

\begin{figure}[tbp]
\centering
\includegraphics[width=0.69\textwidth]{lecture_02_figure_08.png}
\caption{Extraction of a basis component and reconstruction of
\(|v\rangle\). The board preserves the same sequence as
Eqs.~\eqref{eq:l02-component-extraction} and
\eqref{eq:l02-completeness-action}.}
\lectureframe{01:13:45}
\end{figure}

\section{Classical Alternatives and Quantum Superpositions}

We can now attach physical meaning to the vector space. A classical phase
space is a set of possible states. Subsets support the ordinary logic of
union and intersection, OR and AND. Quantum states carry additional linear
structure. Adding two state vectors and multiplying them by complex numbers
can produce another physical state; mere set membership cannot express this.

Begin with a coin. Its classically distinguishable configurations are heads
and tails. We may call them \(H\) and \(T\), or label them \(+1\) and \(-1\),
but those labels are names, not vector lengths. The quantum model assigns
orthonormal states
\begin{equation}
 \langle H|H\rangle=\langle T|T\rangle=1,
 \qquad
 \langle H|T\rangle=0.
 \label{eq:l02-coin-basis}
\end{equation}
The real physical example behind the coin is an electron with two
distinguishable spin outcomes.

The orthogonality has operational content. Configurations that one ideal
measurement can distinguish with certainty are represented by orthogonal
vectors. This is one of the postulates we shall use. For a six-sided die, the
same rule gives six states \(|1\rangle,\ldots,|6\rangle\) satisfying
\begin{equation}
 \langle i|j\rangle=\delta_{ij}.
 \label{eq:l02-die-basis}
\end{equation}

Now comes the step that has no classical analogue:
\begin{equation}
 |C\rangle=\alpha_H|H\rangle+\alpha_T|T\rangle.
 \label{eq:l02-confused-coin}
\end{equation}
This 'confused coin' is a genuine quantum state, neither heads nor tails.
It is a coherent superposition, not a classical statistical mixture whose
member secretly has one definite face. Likewise,
\begin{equation}
 |D\rangle=\sum_{i=1}^{6}\alpha_i|i\rangle
 \label{eq:l02-confused-die}
\end{equation}
is a superposition of the six die outcomes.

\section{Amplitudes, Probabilities, and Normalization}

\begin{classroomqa}
\audiencequestion{How did the coefficients \(\alpha_H\) and \(\alpha_T\)
become probabilities?\lecturetimestamp{01:33:59--01:37:49}}
\lecturerresponse{They do not themselves become probabilities. They are
probability amplitudes, and the Born rule is a postulate:
\[
 P(H)=|\alpha_H|^2,
 \qquad
 P(T)=|\alpha_T|^2.
\]
This is not the whole meaning of the amplitudes; their phases contain more
information than these two probabilities reveal. Operationally, prepare a
large ensemble of systems in the same state, make the same measurement on
each, and identify the probabilities with the limiting frequencies of heads
and tails.}
\end{classroomqa}

If all outcomes are accounted for, their probabilities add to one:
\begin{equation}
 |\alpha_H|^2+|\alpha_T|^2=1.
 \label{eq:l02-coin-normalization}
\end{equation}
Using Eq.~\eqref{eq:l02-coin-basis}, this is precisely
\begin{equation}
 \langle C|C\rangle=1.
 \label{eq:l02-state-normalization}
\end{equation}
For the die,
\begin{equation}
 \sum_{i=1}^{6}|\alpha_i|^2=1,
 \qquad
 \langle D|D\rangle=1.
 \label{eq:l02-die-normalization}
\end{equation}
Physical state vectors are therefore normalized.

\begin{figure}[tbp]
\centering
\includegraphics[width=0.71\textwidth]{lecture_02_figure_09.png}
\caption{Normalization of the two-state superposition written as both a sum
of squared amplitudes and the norm \(\langle C|C\rangle\).}
\lectureframe{01:32:10}
\end{figure}

\begin{classroomqa}
\audiencequestion{For equal heads and tails probabilities, are both
amplitudes \(1/\sqrt{2}\)?\lecturetimestamp{01:38:02--01:39:24}}
\lecturerresponse{Their magnitudes are \(1/\sqrt{2}\), but the amplitudes may
be complex:
\[
 \alpha_H=\frac{e^{i\theta_H}}{\sqrt{2}},
 \qquad
 \alpha_T=\frac{e^{i\theta_T}}{\sqrt{2}}.
\]
Every complex number on the unit circle has unit magnitude. Multiplying such
a number by \(1/\sqrt{2}\) gives probability \(1/2\). For six equally likely
die outcomes, each amplitude has magnitude \(1/\sqrt{6}\).}
\end{classroomqa}

We have not yet said how to prepare an electron in the analogue of this
half-heads, half-tails state. That experimental question will return when
spin replaces the toy coin. For the moment it is more efficient to state the
rules, learn how they work, and then ask which experiments forced physicists
to adopt them. The first rules are now visible:
\begin{enumerate}
\item states are represented by vectors in a complex vector space;
\item perfectly distinguishable alternatives are represented by orthogonal
vectors;
\item physical state vectors have unit norm; and
\item squared amplitude magnitudes give probabilities in an orthonormal
measurement basis.
\end{enumerate}

\section{The First Look at Operators}

There is one more mathematical object to introduce. An operation assigns a
new vector to every vector in the space. Fixed-angle rotation, uniform
rescaling, and reflection about an axis are familiar examples. Squaring a
vector's length is also an operation, but it is not linear: a vector of
length \(2\) is sent to one of length \(4\), while a vector of length \(3\)
is sent to one of length \(9\).

We denote a linear operator by a hatted letter, such as \(\hat L\). It obeys
homogeneity and additivity:
\begin{align}
 \hat L\bigl(\alpha|A\rangle\bigr)
 &=\alpha\,\hat L|A\rangle,\\
 \hat L\bigl(|A\rangle+|B\rangle\bigr)
 &=\hat L|A\rangle+\hat L|B\rangle.
 \label{eq:l02-linearity}
\end{align}
Equivalently, for arbitrary complex \(\alpha,\beta\),
\begin{equation}
 \hat L\bigl(\alpha|A\rangle+\beta|B\rangle\bigr)
 =\alpha\hat L|A\rangle+\beta\hat L|B\rangle.
 \label{eq:l02-combined-linearity}
\end{equation}

\begin{figure}[tbp]
\centering
\includegraphics[width=0.70\textwidth]{lecture_02_figure_10.png}
\caption{The two defining rules for a linear operator: scalar multiplication
passes through the operator, and the image of a sum is the sum of the
images.}
\lectureframe{01:49:20}
\end{figure}

Rotation by a fixed angle passes both tests. Rotating twice a vector gives
twice the rotated vector. If \(\mathbf a+\mathbf b\) is the diagonal of a
parallelogram, rotating the whole parallelogram shows
\begin{equation}
 R(\mathbf a+\mathbf b)=R\mathbf a+R\mathbf b.
\end{equation}
Uniform doubling and reflection are linear for the same reason. The
length-squaring operation fails: scaling the input by \(\alpha\) changes the
length squared by \(|\alpha|^2\), not by \(\alpha\).

\begin{figure}[tbp]
\centering
\begin{tikzpicture}[>=Latex,scale=0.95]
  \coordinate (O) at (0,0);
  \coordinate (A) at (3.4,0.4);
  \coordinate (B) at (1.4,2.5);
  \coordinate (S) at (4.8,2.9);
  \draw[->,thick] (O) -- (A) node[midway,below] {$\mathbf a$};
  \draw[->,thick] (O) -- (B) node[midway,left] {$\mathbf b$};
  \draw[dashed] (A) -- (S);
  \draw[dashed] (B) -- (S);
  \draw[->,very thick,blue!65!black] (O) -- (S)
    node[midway,above left] {$\mathbf a+\mathbf b$};
  \node[align=left] at (7.2,1.45)
    {rotate the whole\\parallelogram};
  \draw[->,bend left=25] (5.0,2.9) to (6.4,2.2);
\end{tikzpicture}
\caption{The parallelogram proof of additivity for a fixed rotation.}
\end{figure}

State vectors describe the possible states of a quantum system. Linear
operators will describe measurable quantities, or observables: the quantum
version of asking whether the coin is heads or tails. The precise relation
between an observable, its operator, and the outcomes of a measurement is the
next step. We are now close enough to the formalism to begin making
predictions.
